
\documentclass{report}

\input{preamble}
\input{macros}
\input{letterfonts}

\title{\Huge{Economia Politica}\\Sistemi Informativi Aziendali}
\author{\huge{Caberlotto Giovanni}}
\date{}

\begin{document}

\maketitle
\newpage% or \cleardoublepage
% \pdfbookmark[<level>]{<title>}{<dest>}
\pdfbookmark[section]{\contentsname}{toc}
\tableofcontents
\pagebreak

\chapter{Economia Pubblica e le Diverse Libertà} % (fold)

\label{chap:Economia Pubblica e le Diverse Libertà}

\section{L'economia pubblica}

\dfn{Economia Pubblica}{L'economia pubblica è la branca della disciplina economica che studia il ruolo del soggetto pubblico (lo stato e più in generale, la pubblica amministrazione) nel sistema economico nonché le possibili forme dell'intervento pubblico, con una impostazione perlo più microeconomica.}

Le scelte pubbliche, soggette al vincolo della scarsità delle risorse, inevitabilmente si traducono in una politica economica che avvantaggia alcune persone svantaggiandone altre. Perciò ogni scelta è il risultato di un'analisi positiva (basata su dati, tabelle e grafici) e di un'analisi normativa (basata su norme di comportamento e giudizi di valore)
In questo modo, l'economia pubblica è in grado di fornire criteri per la scelta di proposte alternative nel campo delle politiche pubbliche.

Alla fine del Medioevo, con l'affermazione degli Stati moderni, si assiste all'organizzazione delle strutture politiche centrali, alla progressiva costituzione di un soggetto pubblico autonomo e alla definizione di nuovi ambiti di intervento pubblico (via via, oltre alla riscossione delle imposte, vengono istituzionalizzate la difesa, la sicurezza, la giustizia, ecc\dots).
L'attenzione dei primi economisti mentre l'economia inizia ad affermarsi come scienza cade sulle nuove forme di intervento dello Stato, e in particolare, sul prelievo fiscale. Per questo motivo l'economia pubblica, soprattutto in Italia, si configura inizialmente come scienza delle finanze.

\dfn{Scienze delle finanze}{In buona sostanza, possiamo dire che la scienza delle finanze analizza e tre funzioni essenziali attribuite all'intervento pubblico (allocativa, redistributiva e di stabilizzazione macroeconomica) in uno specifico contesto istituzionale, che è quello dei sistemi capitalistici}

Non a caso proprio L'italia ha prodotto, tra gli ultimi decenni del $XIX$ secolo e la prima metà del $XX$ secolo, eccellenti studiosi dell'attività finanziaria dello Stato (Einaudi, De Viti De Marco, Mazzola, Cosciani), i quali hanno rivolto il proprio interesse sopratutto sulle entrate tributarie e sulla spesa per l'erogazione di beni e servizi pubblici.
Quando l'intervento pubblico si è spinto ben oltre la sola produzione di beni e servizi destinati alla collettività, intervenendo nella regolamentazione del settore privato e nel funzionamento del mercato, l'ambito di studio della scienza delle finanze si è ampliato al punto che oggi si preferisce parlare di economia pubblica.

Quest'ultima, nel frattempo, ha subito anche l'influenza dell'economia del benessere, che ha impegnato lo Stato nella ricerca di soluzioni a problemi socialmente rilevanti

Dopo la seconda guerra mondiale il campo d'indagine dell'economia pubblica si è organizzato in subdiscipline: economia ambientale, sanitaria, dell'istruzione, dei trasporti ecc.
L'accresciuta presenza dello Stato nell'economia, le molte relazioni tra soggetti pubblici e privati e l'impatto sociale sempre più rilevante dell'intervento pubblico fanno sì che ogni atto pubblico (legge, delibera comunale, sentenza) comporti conseguenze economiche anche per famiglie e imprese.
L'economia pubblica, da sempre fortemente intrecciata con il diritto, in particolare con il diritto tributario, oggi è perciò fortemente intrecciata con numerose scienze sociali, come la scienza politica, la sociologia e la storia economica.


\end{document}
